\documentclass[tikz,border=10pt]{standalone}
\usepackage{tikz}
\usetikzlibrary{calc,arrows.meta,patterns}
\usepackage{amsmath}

\begin{document}

% ==================================================
% Configuration
% ==================================================

% --- Individual vector stack dimensions ---
% Input vectors (embed)
\def\InputEmbedWidth{0.3}
\def\InputEmbedHeight{3.0}
\def\NumInputEmbedVectors{6}

% Input vectors (blend)
\def\InputBlendWidth{0.3}
\def\InputBlendHeight{1.5}
\def\NumInputBlendVectors{6}

% Input vectors (other)
\def\InputOtherWidth{0.3}
\def\InputOtherHeight{0.8}
\def\NumInputOtherVectors{6}

% Output vectors (temporal outputs from encoders)
\def\OutputEmbedWidth{0.3}
\def\OutputEmbedHeight{1.5}

\def\OutputBlendWidth{0.3}
\def\OutputBlendHeight{0.75}

\def\OutputOtherWidth{0.3}
\def\OutputOtherHeight{0.4}

% Final temporal output (inside each block and after pooling)
\def\FinalTemporalWidth{0.3}
\def\FinalTemporalHeight{2.5}

% --- Common vector stacking offsets ---
\def\VectorOffsetX{0.06}
\def\VectorOffsetY{0.06}

% --- Dimension label configuration ---
\def\DimLabelGap{0.3}

% --- Temporal encoder block geometry ---
\def\EncoderWidth{3.2}
\def\EncoderHeight{1.4}
\def\InputToEncoderDist{3.5}
\def\EncoderToOutputDist{2.7}

% --- Fully connected block geometry (inside temporal block) ---
\def\FCWidth{3.2}
\def\FCHeight{1.4}
\def\OutputsToFCDist{4.0}
\def\FCDistToFinalOutput{2.7}

% --- Multiple-instance pooling block geometry ---
\def\MIPoolWidth{3.2}
\def\MIPoolHeight{1.4}
\def\BlocksToPoolDist{5.0}     % distance from front block temporal vector to MI-pooling block
\def\PoolToFinalDist{2.7}      % distance from MI block to pooled temporal vector

% --- Extra FC + outputs after pooling ---
\def\PostPoolFCDist{3.5}         % pooled vector -> post-pool FC center
\def\PostPoolFCToScalarDist{2.5} % post-pool FC -> outputs base (x distance)

% Main scalar output (size 1)
\def\ScalarOutputWidth{0.3}
\def\ScalarOutputHeight{0.3}
\def\ScalarOutputDimText{$1$}

% Auxiliary output (size 9)
\def\AuxOutputWidth{0.3}
\def\AuxOutputHeight{1}
\def\AuxOutputDimText{$9$}
\def\AuxOutputYOffset{2.5}   % vertical offset between scalar and auxiliary outputs

% Show / hide auxiliary head: 1 = show, 0 = hide
\def\ShowAuxOutput{0}

% Style controls for auxiliary loss (striped)
\def\AuxVectorPattern{north east lines}
\def\AuxVectorPatternColor{red!60}

% --- Connector configuration ---
\def\ConnectorGap{0.3}
\def\ConnLineWidth{0.9pt}
\def\ConnColor{black}
\def\ConnStyle{solid}

% --- S-shaped connectors (inside each block) ---
\def\FCConnOutAngle{0}
\def\FCConnInAngle{180}
\def\FCConnLooseness{1.1}

% --- S-shaped connectors (blocks -> MI pooling) ---
\def\PoolConnOutAngle{0}
\def\PoolConnInAngle{180}
\def\PoolConnLooseness{1.1}

% --- Vertical stacking of temporal paths ---
\def\PathVerticalGap{3}

% --- Block repetition configuration ---
\def\NumBlocks{4}              % Number of times to repeat the entire block
\def\BlockOffsetX{-0.2}        % Horizontal offset between blocks
\def\BlockOffsetY{0.2}         % Vertical offset between blocks
\def\BlockBgColor{white}       % Background color for each block
\def\BlockPadding{0.5}         % Padding around the content

% --- Bounding box coordinates (fixed for all blocks) ---
\def\BBoxLeft{-1.5}
\def\BBoxRight{18}
\def\BBoxTop{7}
\def\BBoxBottom{-4}

% --- Text configuration ---
\def\EncoderText{Temporal encoder:\\LSTM}
\def\FCText{Fully connected\\network}
\def\MIPoolText{Multiple instance\\pooling}
\def\TemporalOutputDimText{$d_{\text{temporal}}$}

% --- Styles ---
\tikzset{
  vectorRect/.style={
    draw,
    fill=blue!10,
    rounded corners=1pt
  },
  encoderBlock/.style={
    draw,
    fill=gray!10,
    rounded corners=2pt,
    thick
  },
  fcBlock/.style={
    draw,
    fill=orange!10,
    rounded corners=2pt,
    thick
  },
  miBlock/.style={
    draw,
    fill=green!10,
    rounded corners=2pt,
    thick
  },
  conn/.style={
    -{Latex[length=3mm]},
    line width=\ConnLineWidth,
    draw=\ConnColor,
    \ConnStyle
  }
}

% ==================================================
% Helper macro: draw a stack of 1D vectors (vertical label)
% ==================================================
\newcommand{\drawVectorStack}[6]{%
  \begin{scope}[shift={#6}]
    \pgfmathtruncatemacro{\Last}{#2-1}
    \foreach \i in {0,...,\Last}{%
      \draw[vectorRect]
        (\i * -\VectorOffsetX, \i * \VectorOffsetY)
        rectangle ++(#3, #4);
    }
    \pgfmathsetmacro{\eastX}{\Last * -\VectorOffsetX + #3}
    \pgfmathsetmacro{\eastY}{\Last * \VectorOffsetY + 0.5 * #4}
    \pgfmathsetmacro{\westX}{\Last * -\VectorOffsetX}
    \coordinate (#1-east) at (\eastX,\eastY);
    \coordinate (#1-west) at (\westX,\eastY);
    \node[rotate=90] (#1-dimlabel)
      at ($(#1-west)+(-\DimLabelGap,0)$) {#5};
  \end{scope}
}

% ==================================================
% Helper macro: draw a stack of 1D vectors (horizontal label on west side)
% ==================================================
\newcommand{\drawVectorStackHorizontalLeft}[6]{%
  \begin{scope}[shift={#6}]
    \pgfmathtruncatemacro{\Last}{#2-1}
    \foreach \i in {0,...,\Last}{%
      \draw[vectorRect]
        (\i * -\VectorOffsetX, \i * \VectorOffsetY)
        rectangle ++(#3, #4);
    }
    \pgfmathsetmacro{\eastX}{\Last * -\VectorOffsetX + #3}
    \pgfmathsetmacro{\eastY}{\Last * \VectorOffsetY + 0.5 * #4}
    \pgfmathsetmacro{\westX}{\Last * -\VectorOffsetX}
    \coordinate (#1-east) at (\eastX,\eastY);
    \coordinate (#1-west) at (\westX,\eastY);
    \node (#1-dimlabel)
      at ($(#1-west)+(-\DimLabelGap,0)$) {#5};
  \end{scope}
}

% ==================================================
% Helper macro: one temporal path
% ==================================================
\newcommand{\drawTemporalPath}[9]{%
  \begin{scope}[shift={(0,#2)}]
    \drawVectorStack{#1-in}{#6}{#4}{#5}{#3}{(0,0)}
    \coordinate (#1-encoder-center) at ($(#1-in-east)+(\InputToEncoderDist,0)$);
    \node[encoderBlock,
          minimum width=\EncoderWidth cm,
          minimum height=\EncoderHeight cm,
          align=center]
      (#1-encoder) at (#1-encoder-center) {#7};
    \coordinate (#1-inputConn)   at ($(#1-in-east)+(2*\ConnectorGap,0)$);
    \coordinate (#1-encConnLeft) at ($(#1-encoder.west)+(-\ConnectorGap,0)$);
    \draw[conn] (#1-inputConn) -- (#1-encConnLeft);
    \pgfmathsetmacro{\outputShift}{-0.5*#9}
    \drawVectorStack{#1-out}{1}{\OutputEmbedWidth}{#9}{#8}
      {($(#1-encoder.east)+(\EncoderToOutputDist,\outputShift)$)}
    \coordinate (#1-encConnRight) at ($(#1-encoder.east)+(\ConnectorGap,0)$);
    \coordinate (#1-outConnLeft)  at ($(#1-out-dimlabel.north)+(-\ConnectorGap,0)$);
    \draw[conn] (#1-encConnRight) -- (#1-outConnLeft);
  \end{scope}
}

% ==================================================
% Main content: complete architecture (content only)
% ==================================================
\newcommand{\drawArchitectureContent}{
  \drawTemporalPath{embed}{\PathVerticalGap}%
    {$d_{\text{embed}}$}{\InputEmbedWidth}{\InputEmbedHeight}{\NumInputEmbedVectors}%
    {\EncoderText}{$d_{\text{temporal,embed}}$}{\OutputEmbedHeight};
  \drawTemporalPath{blend}{0}%
    {$d_{\text{blend}}$}{\InputBlendWidth}{\InputBlendHeight}{\NumInputBlendVectors}%
    {\EncoderText}{$d_{\text{temporal,blend}}$}{\OutputBlendHeight};
  \drawTemporalPath{other}{-\PathVerticalGap}%
    {$d_{\text{other}}$}{\InputOtherWidth}{\InputOtherHeight}{\NumInputOtherVectors}%
    {\EncoderText}{$d_{\text{temporal,other}}$}{\OutputOtherHeight};

  \coordinate (fc-center) at ($(blend-out-east)+(\OutputsToFCDist,0)$);
  \node[fcBlock,
        minimum width=\FCWidth cm,
        minimum height=\FCHeight cm,
        align=center]
    (fc) at (fc-center) {\FCText};

  \coordinate (embed-out-conn) at ($(embed-out-east)+(\ConnectorGap,0)$);
  \coordinate (embed-fc-conn)  at ($(fc.west)+(-\ConnectorGap,0)$);
  \draw[conn]
    (embed-out-conn)
      to[out=\FCConnOutAngle, in=\FCConnInAngle, looseness=\FCConnLooseness]
    (embed-fc-conn);

  \coordinate (blend-out-conn) at ($(blend-out-east)+(\ConnectorGap,0)$);
  \coordinate (blend-fc-conn)  at ($(fc.west)+(-\ConnectorGap,0)$);
  \draw[conn]
    (blend-out-conn)
      to[out=\FCConnOutAngle, in=\FCConnInAngle, looseness=\FCConnLooseness]
    (blend-fc-conn);

  \coordinate (other-out-conn) at ($(other-out-east)+(\ConnectorGap,0)$);
  \coordinate (other-fc-conn)  at ($(fc.west)+(-\ConnectorGap,0)$);
  \draw[conn]
    (other-out-conn)
      to[out=\FCConnOutAngle, in=\FCConnInAngle, looseness=\FCConnLooseness]
    (other-fc-conn);

  \pgfmathsetmacro{\finalShift}{-0.5*\FinalTemporalHeight}
  \drawVectorStack{temporal}{1}{\FinalTemporalWidth}{\FinalTemporalHeight}%
    {\TemporalOutputDimText}%
    {($(fc.east)+(\FCDistToFinalOutput,\finalShift)$)}

  \coordinate (fcConnRight)   at ($(fc.east)+(\ConnectorGap,0)$);
  \coordinate (tempConnLeft)  at ($(temporal-dimlabel.north)+(-\ConnectorGap,0)$);
  \draw[conn] (fcConnRight) -- (tempConnLeft);
}

% ==================================================
% Complete block: background + content
% ==================================================
\newcommand{\drawCompleteBlock}{
  \fill[fill=\BlockBgColor, rounded corners=3pt]
    (\BBoxLeft, \BBoxBottom) rectangle (\BBoxRight, \BBoxTop);
  \draw[line width=1pt, rounded corners=3pt]
    (\BBoxLeft, \BBoxBottom) rectangle (\BBoxRight, \BBoxTop);
  \drawArchitectureContent
}

% ==================================================
% Main picture: multiple stacked blocks + MI pooling
% ==================================================
\begin{tikzpicture}

  \pgfmathtruncatemacro{\LastBlock}{\NumBlocks-1}
  \pgfmathsetmacro{\BBoxMidY}{0.5*(\BBoxTop+\BBoxBottom)}

  \foreach \i in {0,...,\LastBlock}{%
    \begin{scope}[shift={(\i*\BlockOffsetX, \i*\BlockOffsetY)}]
      \drawCompleteBlock
      \coordinate (block-\i-east) at (\BBoxRight,\BBoxMidY);
      \coordinate (block-\i-temporal-east) at (temporal-east);
    \end{scope}
  }

  % MI pooling block
  \coordinate (mipool-center) at ($(block-\LastBlock-temporal-east)+(\BlocksToPoolDist,0)$);
  \node[miBlock,
        minimum width=\MIPoolWidth cm,
        minimum height=\MIPoolHeight cm,
        align=center]
    (mipool) at (mipool-center) {\MIPoolText};

  \coordinate (mipool-in) at ($(mipool.west)+(-\ConnectorGap,0)$);
  \foreach \i in {0,...,\LastBlock}{%
    \coordinate (block-\i-out-conn) at (block-\i-east);
    \draw[conn]
      (block-\i-out-conn)
        to[out=\PoolConnOutAngle, in=\PoolConnInAngle, looseness=\PoolConnLooseness]
      (mipool-in);
  }

  % Pooled temporal vector
  \pgfmathsetmacro{\pooledShift}{-0.5*\FinalTemporalHeight}
  \drawVectorStack{temporalPooled}{1}{\FinalTemporalWidth}{\FinalTemporalHeight}%
    {\TemporalOutputDimText}%
    {($(mipool.east)+(\PoolToFinalDist,\pooledShift)$)}

  \coordinate (mipoolConnRight)   at ($(mipool.east)+(\ConnectorGap,0)$);
  \coordinate (tempPooledConnLeft)  at ($(temporalPooled-dimlabel.north)+(-\ConnectorGap,0)$);
  \draw[conn] (mipoolConnRight) -- (tempPooledConnLeft);

  % Post-pooling FC
  \coordinate (postfc-center) at ($(temporalPooled-east)+(\PostPoolFCDist,0)$);
  \node[fcBlock,
        minimum width=\FCWidth cm,
        minimum height=\FCHeight cm,
        align=center]
    (postfc) at (postfc-center) {\FCText};

  \coordinate (tempPooledConnRight) at ($(temporalPooled-east)+(\ConnectorGap,0)$);
  \coordinate (postfcConnLeft)      at ($(postfc.west)+(-\ConnectorGap,0)$);
  \draw[conn]
    (tempPooledConnRight)
      to[out=\FCConnOutAngle, in=\FCConnInAngle, looseness=\FCConnLooseness]
    (postfcConnLeft);

  % Final outputs (size 1 and optional size 9)
  \pgfmathsetmacro{\scalarShift}{-0.5*\ScalarOutputHeight}
  \drawVectorStackHorizontalLeft{scalarOut}{1}{\ScalarOutputWidth}{\ScalarOutputHeight}%
    {\ScalarOutputDimText}%
    {($(postfc.east)+(\PostPoolFCToScalarDist,\scalarShift)$)}

  % Scalar connector: straight horizontal, with gap before label
  \coordinate (postfcConnRight)  at ($(postfc.east)+(\ConnectorGap,0)$);
  \coordinate (scalarConnLeft)   at ($(scalarOut-dimlabel)+(-\ConnectorGap,0)$);
  \draw[conn] (postfcConnRight) -- (scalarConnLeft);

  % Auxiliary head and connector (conditionally shown)
  \ifnum\ShowAuxOutput=1\relax
    \pgfmathsetmacro{\auxShift}{\scalarShift - \AuxOutputYOffset}
    {%
      \tikzset{
        vectorRect/.style={
          draw,
          pattern=\AuxVectorPattern,
          pattern color=\AuxVectorPatternColor,
          rounded corners=1pt
        }
      }
      \drawVectorStackHorizontalLeft{auxOut}{1}{\AuxOutputWidth}{\AuxOutputHeight}%
        {\AuxOutputDimText}%
        {($(postfc.east)+(\PostPoolFCToScalarDist,\auxShift)$)}
    }
    \coordinate (auxConnLeft)      at ($(auxOut-dimlabel)+(-\ConnectorGap,0)$);
    \draw[conn]
      (postfcConnRight)
        to[out=\FCConnOutAngle, in=\FCConnInAngle, looseness=\FCConnLooseness]
      (auxConnLeft);
  \fi

\end{tikzpicture}

\end{document}